% ============================================================
%  يُترجَم بـ XeLaTeX  (في Overleaf: Menu → Compiler → XeLaTeX)
%  الخط: Amiri (متوفر في Overleaf، أو حمّله من github.com/aliftype/amiri)
% ============================================================
\documentclass[12pt,a4paper]{article}
\usepackage[a4paper,margin=1.7cm,top=2.3cm,bottom=2.9cm,headheight=15pt,footskip=1.5cm]{geometry}
\usepackage{amsmath,amssymb}
\usepackage{icomma}
\usepackage[english,bidi=default,layout=lists]{babel}
\babelprovide[import,main]{arabic}
\babelfont{rm}[Scale=1.08]{Amiri}
\babelfont{sf}[Scale=1.08]{Amiri}
\usepackage[most]{tcolorbox}
\usepackage{tikz}
\usetikzlibrary{arrows.meta,calc,shadings}
\usepackage{array}
\usepackage{colortbl}
\usepackage{fancyhdr}
\usepackage{titlesec}
\usepackage{enumitem}
\usepackage{collcell}
\usepackage{needspace}
\usepackage{etoolbox}
\newcolumntype{C}{>{\collectcell\ar}c<{\endcollectcell}}
\pretocmd{\section}{\Needspace{15\baselineskip}}{}{}
\pretocmd{\subsection}{\Needspace{11\baselineskip}}{}{}

\definecolor{cDef}{HTML}{1F4E9C}
\definecolor{cThm}{HTML}{B3261E}
\definecolor{cRem}{HTML}{2E7D32}
\definecolor{cEx}{HTML}{C76B00}
\definecolor{cMeth}{HTML}{6A1B9A}
\definecolor{cSol}{HTML}{00695C}
\definecolor{cSec}{HTML}{1F3A68}

\titleformat{\section}{\Large\bfseries\color{cSec}}{\thesection.}{0.6em}{}[{\color{cSec!60}\titlerule[1pt]}]
\titleformat{\subsection}{\large\bfseries\color{cThm!85!black}}{\thesubsection.}{0.5em}{}
\titlespacing*{\section}{0pt}{14pt}{8pt}
\titlespacing*{\subsection}{0pt}{10pt}{5pt}
\setlength{\parindent}{0pt}
\setlength{\parskip}{4pt}
\setlist{itemsep=2pt,topsep=3pt}

\newtcolorbox{arbox}[3][]{enhanced,
  colback=#2!5,colframe=#2!80!black,arc=2mm,boxrule=0.8pt,
  left=6pt,right=6pt,top=8pt,bottom=5pt,
  fonttitle=\bfseries,coltitle=white,
  title={\foreignlanguage{arabic}{#3}},
  attach boxed title to top right={xshift=-5mm,yshift=-3mm},
  boxed title style={colback=#2!80!black,arc=1.5mm,boxrule=0pt},
  before upper={\selectlanguage{arabic}},
  before skip=8pt,after skip=8pt,breakable,#1}
\newenvironment{BOX}[2]{\begin{otherlanguage}{english}\begin{arbox}{#1}{#2}}{\end{arbox}\end{otherlanguage}}
\newenvironment{tarif}[1][تعريف]{\begin{BOX}{cDef}{#1}}{\end{BOX}}
\newenvironment{khassia}[1][خاصية]{\begin{BOX}{cThm}{#1}}{\end{BOX}}
\newenvironment{molahada}[1][ملاحظة]{\begin{BOX}{cRem}{#1}}{\end{BOX}}
\newenvironment{mithal}[1][مثال]{\begin{BOX}{cEx}{#1}}{\end{BOX}}
\newenvironment{tariqa}[1][طريقة]{\begin{BOX}{cMeth}{#1}}{\end{BOX}}

\newcounter{exo}
\newenvironment{tamrin}[1]{\stepcounter{exo}\begin{BOX}{cSec}{التمرين \theexo: #1}}{\end{BOX}}
\newenvironment{hal}{\Needspace{12\baselineskip}\begin{otherlanguage}{english}\begin{arbox}[breakable]{cSol}{حل التمرين \theexo}}{\end{arbox}\end{otherlanguage}}

\newenvironment{LTR}{\begin{otherlanguage}{english}}{\end{otherlanguage}}
\newcommand{\ar}[1]{\foreignlanguage{arabic}{#1}}

\renewcommand{\dots}{\mathinner{\ldotp\ldotp\ldotp}}
\renewcommand{\ldots}{\mathinner{\ldotp\ldotp\ldotp}}
\renewcommand{\cdots}{\mathinner{\cdotp\cdotp\cdotp}}
\newcommand{\N}{\mathbb{N}}
\newcommand{\Z}{\mathbb{Z}}
\newcommand{\D}{\mathbb{D}}
\newcommand{\Q}{\mathbb{Q}}
\newcommand{\R}{\mathbb{R}}
\newcommand{\pgcd}{\operatorname{PGCD}}
\newcommand{\ppcm}{\operatorname{PPCM}}

% \NLseg{a}{b}{[}{]} ; \NLright{a}{[} ; \NLleft{b}{]}  (interval on a number line)
\newcommand{\NLseg}[4]{\begin{LTR}\begin{tikzpicture}[baseline=-0.6ex]
 \draw[-{Stealth[length=2mm]},gray!70](0,0)--(3.4,0);
 \draw[line width=2pt,cThm!80!black](0.7,0)--(2.7,0);
 \node[cThm!80!black,font=\LARGE] at (0.7,0){$#3$};
 \node[cThm!80!black,font=\LARGE] at (2.7,0){$#4$};
 \node[below=3pt,font=\footnotesize] at (0.7,0){$#1$};
 \node[below=3pt,font=\footnotesize] at (2.7,0){$#2$};
\end{tikzpicture}\end{LTR}}
\newcommand{\NLright}[2]{\begin{LTR}\begin{tikzpicture}[baseline=-0.6ex]
 \draw[-{Stealth[length=2mm]},gray!70](0,0)--(3.4,0);
 \draw[line width=2pt,cThm!80!black](1.3,0)--(3.3,0);
 \node[cThm!80!black,font=\LARGE] at (1.3,0){$#2$};
 \node[below=3pt,font=\footnotesize] at (1.3,0){$#1$};
\end{tikzpicture}\end{LTR}}
\newcommand{\NLleft}[2]{\begin{LTR}\begin{tikzpicture}[baseline=-0.6ex]
 \draw[-{Stealth[length=2mm]},gray!70](0,0)--(3.4,0);
 \draw[line width=2pt,cThm!80!black](0.05,0)--(2.1,0);
 \node[cThm!80!black,font=\LARGE] at (2.1,0){$#2$};
 \node[below=3pt,font=\footnotesize] at (2.1,0){$#1$};
\end{tikzpicture}\end{LTR}}

% ---------- social-media icons (pure TikZ) ----------
\definecolor{cFB}{HTML}{1877F2}
\definecolor{cTG}{HTML}{229ED9}
\definecolor{cIG1}{HTML}{FEDA75}
\definecolor{cIG2}{HTML}{D62976}
\definecolor{cIG3}{HTML}{4F5BD5}
\newcommand{\fbicon}{\begin{LTR}\begin{tikzpicture}[baseline=-0.75ex,x=1em,y=1em]
 \fill[cFB] (0,0) circle (0.55);
 \node[white,inner sep=0,font=\usefont{TU}{lmss}{bx}{n}\small] at (0.04,-0.02) {f};
\end{tikzpicture}\end{LTR}}
\newcommand{\igicon}{\begin{LTR}\begin{tikzpicture}[baseline=-0.75ex,x=1em,y=1em]
 \shade[lower left=cIG1,lower right=cIG2,upper left=cIG2,upper right=cIG3,rounded corners=0.28em] (-0.55,-0.55) rectangle (0.55,0.55);
 \draw[white,line width=0.09em,rounded corners=0.18em] (-0.36,-0.36) rectangle (0.36,0.36);
 \draw[white,line width=0.09em] (0,0) circle (0.17);
 \fill[white] (0.22,0.22) circle (0.05);
\end{tikzpicture}\end{LTR}}
\newcommand{\tgicon}{\begin{LTR}\begin{tikzpicture}[baseline=-0.75ex,x=1em,y=1em]
 \fill[cTG] (0,0) circle (0.55);
 \fill[white] (-0.33,0.0) -- (0.32,0.27) -- (0.2,-0.3) -- (0.0,-0.15) -- (-0.08,-0.28) -- (-0.08,-0.08) -- cycle;
 \fill[cTG!70!black] (-0.08,-0.08) -- (0.25,0.2) -- (0.0,-0.15) -- (-0.08,-0.28) -- cycle;
\end{tikzpicture}\end{LTR}}
% page number: black box, white two-digit number, at the right of the social row
\newcommand{\PageBox}{\begin{LTR}\tikz[baseline=(p.base)]\node[fill=black,text=white,inner sep=2.5pt,minimum width=1.6em,minimum height=1.3em,font=\usefont{TU}{lmss}{bx}{n}\small](p){\ifnum\value{page}<10 0\fi\arabic{page}};\end{LTR}}
% visual left-to-right order: Telegram | Instagram | Facebook | page box
\newcommand{\SocialRow}{\begin{otherlanguage}{english}%
 \ar{\textbf{الأستاذ سوايسية هشام للرياضيات}}\ \tgicon\hfill
 \textbf{prof\_hicham\_math}\ \igicon\hfill
 \ar{\textbf{الأستاذ سوايسية هشام للرياضيات}}\ \fbicon\hspace{2.2em}\PageBox\end{otherlanguage}}

\pagestyle{fancy}
\fancyhf{}
\renewcommand{\headrulewidth}{0.6pt}
\renewcommand{\footrulewidth}{0.4pt}
% \SetHeader{right head: document title}{left head: level -- stream}
\newcommand{\SetHeader}[2]{%
 \fancyhead[C]{\parbox{\textwidth}{\ar{#1\hfill #2}}}%
 \fancyfoot[C]{\parbox{\textwidth}{\small\SocialRow}}}
\usepackage{tkz-tab}
\usepackage{multicol}
% variation / sign table (LTR)
\newenvironment{TAB}{\par\begin{center}\begin{LTR}\begin{tikzpicture}[scale=0.95,every node/.style={font=\small}]}{\end{tikzpicture}\end{LTR}\end{center}}
% grid + axes, clipped: \begin{GR}{xmin}{xmax}{ymin}{ymax}{unit cm}
\newenvironment{GR}[5]{\par\begin{center}\begin{LTR}\begin{tikzpicture}[x=#5cm,y=#5cm,font=\footnotesize]
 \draw[gray!25,very thin] (#1,#3) grid (#2,#4);
 \draw[-{Stealth[length=2mm]}] (#1,0)--(#2+0.3,0) node[below left]{$x$};
 \draw[-{Stealth[length=2mm]}] (0,#3)--(0,#4+0.3) node[below left]{$y$};
 \node[below left] at (0,0){$O$};
 \begin{scope}\clip (#1,#3) rectangle (#2,#4);}
{\end{scope}\end{tikzpicture}\end{LTR}\end{center}}
\newcommand{\Ticks}[2]{\foreach \t in {#1} \node[below,font=\scriptsize] at (\t,0){$\t$}; \foreach \t in {#2} \node[left,font=\scriptsize] at (0,\t){$\t$};}
\newcommand{\bx}[1]{\ensuremath{\colorbox{cDef!12}{$\displaystyle #1$}}}
\newcommand{\TwoCol}[4]{\par\begin{LTR}\noindent\begin{minipage}[t]{#1\textwidth}\selectlanguage{arabic}#2\end{minipage}\hfill\begin{minipage}[t]{#3\textwidth}\selectlanguage{arabic}#4\end{minipage}\end{LTR}\par}
\newcommand{\RL}[4]{\TwoCol{#3}{#4}{#1}{#2}}  % \RL{w1}{right content}{w2}{left content}
\newcommand{\ok}{$\textcolor{cRem}{\checkmark}$}
\newcommand{\limn}{\displaystyle\lim_{n\to+\infty}}
% sub-questions, each on its own line
\newenvironment{sq}{\begin{itemize}[leftmargin=2.4em,itemsep=3pt]}{\end{itemize}}
\newcommand{\qa}{\item[(أ)]}\newcommand{\qb}{\item[(ب)]}\newcommand{\qc}{\item[(جـ)]}\newcommand{\qd}{\item[(د)]}\newcommand{\qe}{\item[(هـ)]}
\newcommand{\pa}{\par\smallskip\textbf{(أ)} }\newcommand{\pb}{\par\smallskip\textbf{(ب)} }\newcommand{\pc}{\par\smallskip\textbf{(جـ)} }\newcommand{\pd}{\par\smallskip\textbf{(د)} }\newcommand{\pe}{\par\smallskip\textbf{(هـ)} }
\newcommand{\QQ}[1]{\par\smallskip\textbf{#1.} }
\newenvironment{tmr}{\stepcounter{exo}\begin{BOX}{cSec}{تمرين \theexo}}{\end{BOX}}
\newcommand{\ABC}[3]{\\[2pt]\begin{tabular}{p{0.28\textwidth}p{0.28\textwidth}p{0.28\textwidth}}(جـ) \ #3 & (ب) \ #2 & (أ) \ #1\end{tabular}}
\pretocmd{\section}{\Needspace{24\baselineskip}}{}{}
\newtcbox{\bxtbox}{on line,colback=cDef!12,colframe=cDef!12,boxrule=0pt,arc=1pt,boxsep=0pt,left=2pt,right=2pt,top=1pt,bottom=1pt}
\newcommand{\bxt}[1]{\foreignlanguage{english}{\bxtbox{\ar{#1}}}}
\usepackage{pgfplots}
\pgfplotsset{compat=1.16}
\pgfplotsset{sc/.style={width=0.8\textwidth,height=6.3cm,axis lines=left,grid=both,minor tick num=1,
  major grid style={gray!40},minor grid style={gray!15},tick label style={font=\scriptsize},label style={font=\footnotesize},
  /pgf/number format/use comma,/pgf/number format/1000 sep={},scaled ticks=false,
  xlabel style={at={(ticklabel* cs:1)},anchor=north east},ylabel style={rotate=-90,at={(ticklabel* cs:1)},anchor=south west}},
  pt/.style={only marks,mark=*,mark size=2.2pt,color=cDef},
  gp/.style={only marks,mark=square*,mark size=2.6pt,color=cThm},
  gq/.style={only marks,mark=triangle*,mark size=3pt,color=cEx},
  ln/.style={thick,color=cThm,samples=2},
  lm/.style={thick,dashed,color=cRem,samples=2}}
% scatter plot on graph paper: \begin{SC}{axis options}
\newenvironment{SC}[1]{\par\begin{center}\begin{LTR}\begin{tikzpicture}\begin{axis}[sc,#1]}{\end{axis}\end{tikzpicture}\end{LTR}\end{center}}
\SetHeader{\small\bfseries\color{cSec}تمارين تجريبية: الإحصاء}{\small\color{cSec}السنة الثالثة ثانوي -- تسيير واقتصاد}
\BeforeBeginEnvironment{tamrin}{\Needspace{12\baselineskip}}
\BeforeBeginEnvironment{hal}{\Needspace{20\baselineskip}}
\begin{document}
\begin{LTR}
\begin{tcolorbox}[enhanced,colback=cSec,colframe=cSec,arc=3mm,halign=center,boxrule=0pt,top=12pt,bottom=12pt,before upper={\selectlanguage{arabic}}]
{\color{white}\fontsize{24}{32}\selectfont\bfseries تمارين تجريبية: الإحصاء}\\[8pt]
{\color{white}\large سحابة النقط -- طريقة ماير -- تعديل بياني ومستقيم معطى -- تغيير المتغير}\\[5pt]
{\color{white}\large 11 تمرينا متدرجا -- بدون حلول (للتدرّب)}\\[5pt]{\color{white}\large السنة الثالثة ثانوي -- تسيير واقتصاد}
\end{tcolorbox}
\end{LTR}

\begin{tarif}[تذكير: القوانين التي نستعملها]
$(x_i;y_i)$ سلسلة إحصائية ذات متغيرين عدد نقطها $n$.
\begin{itemize}
\item \textbf{النقطة المتوسطة:} $G(\bar x;\bar y)$ حيث $\bar x=\dfrac{1}{n}\sum x_i$ و $\bar y=\dfrac{1}{n}\sum y_i$.
\item \textbf{تعديل بياني:} نرسم مستقيما يشمل $G$ ويترك تقريبا نفس عدد النقط من كل جهة، ثم نعيّن معادلته من $G$ ونقطة ثانية منه.
\item \textbf{مستقيم تعديل معطى:} نتحقق أنه يشمل $G$ أي $\bar y=a\bar x+b$ (بتقريب التدوير)، ثم نرسمه بنقطتين.
\item \textbf{طريقة ماير:} نرتب النقط حسب $x_i$، نقسمها إلى مجموعتين متساويتين في العدد، نحسب النقطتين المتوسطتين $G_1$ و $G_2$؛ المستقيم $(G_1G_2)$ هو مستقيم ماير ويشمل $G$.
\item \textbf{تغيير المتغير:} إذا كان $z=\ln y$ و $z=ax+b$ فإن $y=e^{ax+b}=e^{b}\left(e^{a}\right)^{x}$، أي $y=k\,q^{x}$ مع $k=e^{b}$ و $q=e^{a}$.
\end{itemize}
\end{tarif}

\section{سحابة النقط، النقطة المتوسطة وطريقة ماير}
% ------------------------------------------------------------------
\begin{tamrin}{مبيعات متجر -- تعديل بطريقة ماير}
يعطي الجدول التالي المبيعات الشهرية $y_i$ (بالألف دينار جزائري) لمتجر خلال الأشهر الستة الأولى من سنة 2025، حيث $x_i$ رتبة الشهر ($x_i=1$ يوافق جانفي).
\begin{center}\begin{LTR}\renewcommand{\arraystretch}{1.35}\small
\begin{tabular}{|c|c|c|c|c|c|c|}\hline
\rowcolor{cSec!12}\ar{رتبة الشهر} $x_i$ & $1$ & $2$ & $3$ & $4$ & $5$ & $6$\\\hline
\ar{المبيعات} $y_i$ & $420$ & $470$ & $490$ & $520$ & $590$ & $630$\\\hline
\end{tabular}\end{LTR}\end{center}
\QQ{1} مثّل سحابة النقط $M_i(x_i;y_i)$ في معلم متعامد (على محور الفواصل: $1\,\text{cm}$ لكل شهر؛ وعلى محور التراتيب: $1\,\text{cm}$ لكل $20$ ألف دينار، ونبدأ التدريج من $400$).
\QQ{2} احسب إحداثيي النقطة المتوسطة $G$ لهذه السلسلة ثم علّمها.
\QQ{3} هل يبرّر شكل السحابة إجراء تعديل خطي؟
\QQ{4} نقسم السلسلة إلى مجموعتين: الأولى تضم النقط الثلاث الأولى، والثانية تضم النقط الثلاث الأخيرة.
\begin{sq}
\qa احسب إحداثيي $G_1$ و $G_2$ النقطتين المتوسطتين للمجموعتين.
\qb عيّن معادلة مستقيم ماير $(G_1G_2)$ وتحقق أن $G$ تنتمي إليه، ثم ارسمه.
\end{sq}
\QQ{5} باستعمال هذا التعديل:
\begin{sq}
\qa قدّر مبيعات شهر سبتمبر 2025.
\qb ابتداء من أي شهر من سنة 2025 تتجاوز المبيعات $800$ ألف دينار (إذا استمر الاتجاه نفسه)؟
\end{sq}
\end{tamrin}

% ------------------------------------------------------------------
\begin{tamrin}{السعر والطلب -- مستقيم ماير والإيراد الأعظمي}
درست مؤسسة الطلب اليومي $y_i$ (عدد الوحدات المطلوبة) على منتوج بدلالة سعر الوحدة $x_i$ (بالدينار) فتحصلت على:
\begin{center}\begin{LTR}\renewcommand{\arraystretch}{1.35}\small
\begin{tabular}{|c|c|c|c|c|c|c|}\hline
\rowcolor{cSec!12}\ar{السعر} $x_i$ & $20$ & $30$ & $40$ & $50$ & $60$ & $70$\\\hline
\ar{الطلب} $y_i$ & $412$ & $381$ & $350$ & $302$ & $275$ & $251$\\\hline
\end{tabular}\end{LTR}\end{center}
\QQ{1} مثّل سحابة النقط $M_i(x_i;y_i)$ (على محور الفواصل: $1\,\text{cm}$ لكل $10$ دنانير؛ وعلى محور التراتيب: $1\,\text{cm}$ لكل $20$ وحدة ابتداء من $200$).
\QQ{2} احسب إحداثيي النقطة المتوسطة $G$.
\QQ{3} باستعمال طريقة ماير (مجموعة أولى: النقط الثلاث الأولى) بيّن أن معادلة مستقيم التعديل هي: $y=-3{,}5x+486$.
\QQ{4} فسّر اقتصاديا المعامل $-3{,}5$.
\QQ{5}
\begin{sq}
\qa قدّر الطلب اليومي إذا كان سعر الوحدة $80$ دينارا.
\qb ما السعر الذي يوافق طلبا يوميا قدره $300$ وحدة؟ (دوّر إلى $10^{-2}$).
\end{sq}
\QQ{6} نعتبر أن الطلب معطى بـ $y=-3{,}5x+486$ من أجل $x\in[20;100]$. الإيراد اليومي هو $R(x)=x\cdot y$.
\begin{sq}
\qa بيّن أن $R(x)=-3{,}5x^{2}+486x$.
\qb ادرس اتجاه تغير الدالة $R$ على $[20;100]$، واستنتج السعر الذي يجعل الإيراد أعظميا وقيمة هذا الإيراد (دوّر إلى $10^{-2}$).
\end{sq}
\end{tamrin}

% ------------------------------------------------------------------
\begin{tamrin}{تطور عدد سكان مدينة -- ماير بثماني نقط}
يمثل الجدول التالي عدد سكان مدينة $y_i$ (بالآلاف) بين سنتي 2016 و 2023، حيث $x_i$ رتبة السنة.
\begin{center}\begin{LTR}\renewcommand{\arraystretch}{1.35}\small
\begin{tabular}{|c|c|c|c|c|c|c|c|c|}\hline
\rowcolor{cSec!12}\ar{السنة} & $2016$ & $2017$ & $2018$ & $2019$ & $2020$ & $2021$ & $2022$ & $2023$\\\hline
\ar{الرتبة} $x_i$ & $1$ & $2$ & $3$ & $4$ & $5$ & $6$ & $7$ & $8$\\\hline
\ar{السكان} $y_i$ & $24,2$ & $24,9$ & $25,7$ & $26,4$ & $27,5$ & $28,6$ & $29,4$ & $30,1$\\\hline
\end{tabular}\end{LTR}\end{center}
\QQ{1} مثّل سحابة النقط $M_i(x_i;y_i)$ (على محور الفواصل: $1\,\text{cm}$ لكل سنة؛ وعلى محور التراتيب: $1\,\text{cm}$ لكل ألف نسمة ابتداء من $24$).
\QQ{2} احسب إحداثيي النقطة المتوسطة $G$.
\QQ{3} عيّن معادلة مستقيم ماير $(G_1G_2)$، حيث $G_1$ و $G_2$ النقطتان المتوسطتان للنقط الأربع الأولى وللنقط الأربع الأخيرة على الترتيب، ثم تحقق أن $G\in(G_1G_2)$.
\QQ{4} فسّر المعامل الموجه لهذا المستقيم.
\QQ{5}
\begin{sq}
\qa قدّر عدد سكان المدينة سنة 2026.
\qb في أي سنة يتجاوز عدد السكان $35$ ألف نسمة لأول مرة؟
\end{sq}
\QQ{6} احسب النسبة المئوية لزيادة عدد السكان بين سنتي 2016 و 2023 (دوّر إلى $10^{-2}$).
\end{tamrin}

\section{مستقيم تعديل معطى، تعديل بياني، وماير في سياقات اقتصادية}
% ------------------------------------------------------------------
\begin{tamrin}{تكلفة الإنتاج -- مستقيم تعديل معطى}
يعطي الجدول التالي التكلفة الإجمالية $y_i$ (بالألف دينار) لإنتاج $x_i$ مئة وحدة من سلعة في ورشة:
\begin{center}\begin{LTR}\renewcommand{\arraystretch}{1.35}\small
\begin{tabular}{|c|c|c|c|c|c|}\hline
\rowcolor{cSec!12}\ar{الإنتاج} $x_i$ & $1$ & $2$ & $3$ & $4$ & $5$\\\hline
\ar{التكلفة} $y_i$ & $120$ & $150$ & $190$ & $230$ & $260$\\\hline
\end{tabular}\end{LTR}\end{center}
\QQ{1} مثّل سحابة النقط $M_i(x_i;y_i)$ (على محور الفواصل: $2\,\text{cm}$ لكل مئة وحدة؛ وعلى محور التراتيب: $1\,\text{cm}$ لكل $20$ ألف دينار ابتداء من $80$)، ثم احسب إحداثيي النقطة المتوسطة $G$.
\QQ{2} هل يبرّر شكل السحابة تعديلا خطيا؟
\QQ{3} نقبل أن المستقيم $(D):\ y=36x+82$ مستقيم تعديل لهذه السحابة. تحقق أن $G\in(D)$ ثم ارسم $(D)$.
\QQ{4} فسّر اقتصاديا العددين $36$ و $82$.
\QQ{5}
\begin{sq}
\qa قدّر تكلفة إنتاج $700$ وحدة.
\qb ميزانية الورشة $400$ ألف دينار. ما أكبر عدد من الوحدات يمكن إنتاجه؟
\end{sq}
\end{tamrin}

% ------------------------------------------------------------------
\begin{tamrin}{الإشهار ورقم الأعمال -- ماير بقيم غير منتظمة}
يعطي الجدول التالي مصاريف الإشهار $x_i$ (بعشرات الآلاف من الدنانير) ورقم الأعمال $y_i$ (بملايين الدنانير) لمؤسسة خلال ستة أشهر:
\begin{center}\begin{LTR}\renewcommand{\arraystretch}{1.35}\small
\begin{tabular}{|c|c|c|c|c|c|c|}\hline
\rowcolor{cSec!12}\ar{الإشهار} $x_i$ & $2$ & $3$ & $5$ & $6$ & $8$ & $12$\\\hline
\ar{رقم الأعمال} $y_i$ & $2,9$ & $3,8$ & $4,6$ & $5,4$ & $6,1$ & $8,6$\\\hline
\end{tabular}\end{LTR}\end{center}
\QQ{1} مثّل سحابة النقط $M_i(x_i;y_i)$ في معلم متعامد، واحسب إحداثيي $G$.
\QQ{2} $G_1$ و $G_2$ النقطتان المتوسطتان للنقط الثلاث الأولى وللنقط الثلاث الأخيرة.
\begin{sq}
\qa احسب إحداثيي $G_1$ و $G_2$.
\qb بيّن أن معادلة مستقيم ماير $(G_1G_2)$ هي $y=0{,}55x+\dfrac{29}{15}$، وتحقق أنه يشمل $G$.
\end{sq}
\QQ{3}
\begin{sq}
\qa قدّر رقم الأعمال إذا خصصت المؤسسة $150\,000$ دينار للإشهار.
\qb ما مبلغ الإشهار اللازم لبلوغ رقم أعمال قدره $10$ ملايين دينار؟
\end{sq}
\end{tamrin}

% ------------------------------------------------------------------
\begin{tamrin}{مشتركو الإنترنت -- تعديل بياني يمر بـ $G$}
يعطي الجدول عدد المشتركين $y_i$ (بالآلاف) في خدمة الإنترنت عبر الألياف البصرية في بلدية بين 2018 و 2024:
\begin{center}\begin{LTR}\renewcommand{\arraystretch}{1.35}\small
\begin{tabular}{|c|c|c|c|c|c|c|c|}\hline
\rowcolor{cSec!12}\ar{السنة} & $2018$ & $2019$ & $2020$ & $2021$ & $2022$ & $2023$ & $2024$\\\hline
\ar{المشتركون} $y_i$ & $3,2$ & $3,9$ & $4,5$ & $5,4$ & $6$ & $6,8$ & $7,5$\\\hline
\end{tabular}\end{LTR}\end{center}
\QQ{1} نضع $x_i=a_i-2018$ حيث $a_i$ السنة. أعط قيم $x_i$ ثم مثّل سحابة النقط $M_i(x_i;y_i)$.
\QQ{2} احسب إحداثيي النقطة المتوسطة $G$ (دوّر إلى $10^{-4}$).
\QQ{3} نعدّل السحابة بيانيا بالمستقيم $(\Delta)$ الذي يشمل $G$ والنقطة $M_1(0;3{,}2)$.
\begin{sq}
\qa بيّن أن $(\Delta)$ يترك نفس عدد النقط من كل جهة.
\qb عيّن معادلة $(\Delta)$ (دوّر المعامل الموجه إلى $10^{-2}$).
\end{sq}
\QQ{4}
\begin{sq}
\qa قدّر عدد المشتركين سنة 2027.
\qb ابتداء من أي سنة يتجاوز عدد المشتركين $10$ آلاف؟
\end{sq}
\end{tamrin}

% ------------------------------------------------------------------
\begin{tamrin}{العرض والطلب -- سعر التوازن}
في سوق الجملة، سجّل عدد من المنتجين الكمية المعروضة $y_i$ (بالأطنان) من البطاطا بدلالة سعر الكيلوغرام $x_i$ (بالدينار):
\begin{center}\begin{LTR}\renewcommand{\arraystretch}{1.35}\small
\begin{tabular}{|c|c|c|c|c|c|c|}\hline
\rowcolor{cSec!12}\ar{السعر} $x_i$ & $40$ & $45$ & $50$ & $55$ & $60$ & $65$\\\hline
\ar{الكمية المعروضة} $y_i$ & $120$ & $131$ & $139$ & $152$ & $160$ & $171$\\\hline
\end{tabular}\end{LTR}\end{center}
\QQ{1} مثّل سحابة النقط $M_i(x_i;y_i)$ (المبدأ $(35;110)$)، واحسب إحداثيي $G$.
\QQ{2} باستعمال طريقة ماير (النقط الثلاث الأولى ثم الثلاث الأخيرة)، بيّن أن معادلة مستقيم العرض $(D_1)$ هي $y=2{,}07x+37$ (المعامل الموجه مدوّر إلى $10^{-2}$).
\QQ{3} دراسة أخرى بيّنت أن الكمية المطلوبة تعطى بالمستقيم $(D_2):\ y=-1{,}93x+277$.
\begin{sq}
\qa عيّن سعر التوازن (السعر الذي تتساوى عنده الكميتان المعروضة والمطلوبة) والكمية الموافقة.
\qb إذا حُدّد السعر بـ $65$ DA، احسب الفائض في العرض.
\end{sq}
\end{tamrin}

\section{التعديل بتغيير المتغير $z=\ln y$}
% ------------------------------------------------------------------
\begin{tamrin}{طلبيات متجر إلكتروني -- نمو أسي}
يعطي الجدول عدد الطلبيات الشهرية $y_i$ لمتجر إلكتروني خلال الأشهر الستة الأولى من نشاطه ($x_i$ رتبة الشهر):
\begin{center}\begin{LTR}\renewcommand{\arraystretch}{1.35}\small
\begin{tabular}{|c|c|c|c|c|c|c|}\hline
\rowcolor{cSec!12}\ar{الشهر} $x_i$ & $1$ & $2$ & $3$ & $4$ & $5$ & $6$\\\hline
\ar{الطلبيات} $y_i$ & $120$ & $165$ & $230$ & $310$ & $430$ & $590$\\\hline
\end{tabular}\end{LTR}\end{center}
\QQ{1} مثّل سحابة النقط $M_i(x_i;y_i)$. هل شكل السحابة يوحي بتعديل خطي؟
\QQ{2} نضع $z=\ln y$. انقل ثم أتمم الجدول بقيم $z_i$ مدوّرة إلى $10^{-2}$.
\QQ{3}
\begin{sq}
\qa مثّل سحابة النقط $N_i(x_i;z_i)$.
\qb باستعمال طريقة ماير (النقط الثلاث الأولى ثم الثلاث الأخيرة)، عيّن معادلة مستقيم تعديل $z$ بدلالة $x$ (دوّر $a$ و $b$ إلى $10^{-3}$).
\end{sq}
\QQ{4} استنتج أن $y\approx k\,q^{x}$ حيث يطلب تعيين $k$ (إلى $10^{-2}$) و $q$ (إلى $10^{-3}$)، ثم فسّر $q$.
\QQ{5}
\begin{sq}
\qa قدّر عدد الطلبيات في الشهر التاسع.
\qb ابتداء من أي شهر يتجاوز عدد الطلبيات $2000$ طلبية؟
\end{sq}
\end{tamrin}

% ------------------------------------------------------------------
\begin{tamrin}{سكان مدينة جديدة -- ماير بعد تغيير المتغير}
يعطي الجدول عدد سكان $y_i$ (بالآلاف) لمدينة جديدة بين 2017 و 2024:
\begin{center}\begin{LTR}\renewcommand{\arraystretch}{1.35}\small
\begin{tabular}{|c|c|c|c|c|c|c|c|c|}\hline
\rowcolor{cSec!12}\ar{السنة} & $2017$ & $2018$ & $2019$ & $2020$ & $2021$ & $2022$ & $2023$ & $2024$\\\hline
\ar{الرتبة} $x_i$ & $1$ & $2$ & $3$ & $4$ & $5$ & $6$ & $7$ & $8$\\\hline
\ar{السكان} $y_i$ & $50$ & $53,1$ & $56$ & $59,4$ & $63$ & $66,6$ & $70,6$ & $74,9$\\\hline
\end{tabular}\end{LTR}\end{center}
\QQ{1} نضع $z=\ln y$. أعط قيم $z_i$ مدوّرة إلى $10^{-2}$.
\QQ{2} مثّل سحابة النقط $N_i(x_i;z_i)$. هل يبرّر شكلها تعديلا خطيا؟
\QQ{3} باستعمال طريقة ماير (النقط الأربع الأولى ثم الأربع الأخيرة)، بيّن أن معادلة المستقيم المعدِّل هي \mbox{$z=0{,}058x+3{,}852$} (معاملات مدوّرة إلى $10^{-3}$).
\QQ{4} استنتج عبارة $y$ بدلالة $x$ على الشكل $y=k\,q^{x}$ ، ثم استنتج نسبة النمو السنوي للسكان.
\QQ{5}
\begin{sq}
\qa قدّر عدد سكان المدينة سنة 2028.
\qb في أي سنة يتجاوز عدد السكان $100$ ألف نسمة؟
\end{sq}
\end{tamrin}

% ------------------------------------------------------------------
\begin{tamrin}{اهتلاك سيارة -- نقصان أسي}
يعطي الجدول قيمة سيارة $y_i$ (بملايين الدنانير) بدلالة عمرها $x_i$ (بالسنوات) في سوق السيارات المستعملة:
\begin{center}\begin{LTR}\renewcommand{\arraystretch}{1.35}\small
\begin{tabular}{|c|c|c|c|c|c|c|c|}\hline
\rowcolor{cSec!12}\ar{العمر} $x_i$ & $0$ & $1$ & $2$ & $3$ & $4$ & $5$ & $6$\\\hline
\ar{القيمة} $y_i$ & $2,4$ & $2,04$ & $1,75$ & $1,49$ & $1,27$ & $1,08$ & $0,93$\\\hline
\end{tabular}\end{LTR}\end{center}
\QQ{1} مثّل سحابة النقط $M_i(x_i;y_i)$ وصف شكلها.
\QQ{2} نضع $z=\ln y$. أتمم جدول القيم بـ $z_i$ مدوّرة إلى $10^{-2}$، ثم احسب إحداثيي النقطة المتوسطة $G'$ للسلسلة $(x_i;z_i)$.
\QQ{3} مثّل السحابة $N_i(x_i;z_i)$. أي التعديلين أنسب: $y$ بدلالة $x$ أم $z$ بدلالة $x$؟
\QQ{4} نقبل أن المستقيم $(\Delta):\ z=-0{,}158x+0{,}875$ مستقيم تعديل لـ $z$ بدلالة $x$. تحقق أنه يشمل $G'$ (بتقريب $10^{-2}$).
\QQ{5} استنتج أن $y\approx2{,}40\times(0{,}854)^{x}$، ثم استنتج نسبة الاهتلاك السنوي.
\QQ{6}
\begin{sq}
\qa قدّر قيمة السيارة بعد $8$ سنوات.
\qb ابتداء من أي عمر تصبح قيمة السيارة أقل من نصف قيمتها وهي جديدة ($2{,}40$ مليون دينار)؟
\end{sq}
\end{tamrin}

\section{تمرين تقييمي}
% ------------------------------------------------------------------
\begin{tamrin}{اختيار من متعدد وصحيح أم خاطئ}
نعتبر السلسلة الإحصائية ذات المتغيرين:
\begin{center}\begin{LTR}\renewcommand{\arraystretch}{1.35}\small
\begin{tabular}{|c|c|c|c|c|c|c|}\hline
\rowcolor{cSec!12}$x_i$ & $1$ & $2$ & $3$ & $4$ & $5$ & $6$\\\hline
$y_i$ & $3$ & $5$ & $4$ & $7$ & $6$ & $8$\\\hline
\end{tabular}\end{LTR}\end{center}
\textbf{الجزء الأول:} اختر الجواب الصحيح الوحيد في كل حالة.
\QQ{1} النقطة المتوسطة $G$ هي: \ABC{$G(3{,}5;5{,}5)$}{$G(5{,}5;3{,}5)$}{$G(3;5{,}5)$}
\QQ{2} النقطة المتوسطة $G_1$ للنقط الثلاث الأولى هي: \ABC{$G_1(2;4)$}{$G_1(2;5)$}{$G_1(4;2)$}
\QQ{3} مستقيم ماير $(G_1G_2)$ (النقط الثلاث الأولى ثم الثلاث الأخيرة) هو: \ABC{$y=x+2$}{$y=2x+1$}{$y=x+3{,}5$}
\QQ{4} حسب مستقيم ماير، تبلغ القيمة المقدَّرة لـ $y$ العدد $20$ لما: \ABC{$x=18$}{$x=22$}{$x=10$}

\textbf{الجزء الثاني:} أجب بصحيح أو خاطئ مع التعليل.
\QQ{5} مستقيم ماير $(G_1G_2)$ يشمل النقطة المتوسطة $G$.
\QQ{6} المستقيم $(\Delta):\ y=1{,}2x+1{,}5$ يشمل النقطة $G$.
\QQ{7} إذا أضفنا $10$ إلى كل القيم $y_i$ فإن $G$ تصبح $G(3{,}5;15{,}5)$ ومعامل توجيه مستقيم ماير لا يتغير.
\QQ{8} القيمة المقدَّرة لـ $y$ من أجل $x=10$ بمستقيم ماير هي $13$.
\end{tamrin}
\end{document}
